\documentclass[12pt]{article}

\usepackage[utf8]{inputenc}
\usepackage{amsmath, amssymb, amsthm}
\usepackage{geometry}
\geometry{a4paper, margin=2.5cm}

\title{Teorema lui Fréchet}
\author{}
\date{}

\theoremstyle{definition}
\newtheorem{definition}{Definiție}

\theoremstyle{plain}
\newtheorem{theorem}{Teoremă}
\newtheorem{proposition}{Propoziție}

\begin{document}

\maketitle

\begin{definition}
Fie $X$ și $Y$ două spații Banach. Un operator liniar $T : X \to Y$ se numește \emph{închis} dacă pentru orice șir $(x_n)$ din $X$ cu


\[
x_n \to x \quad \text{în } X, \qquad Tx_n \to y \quad \text{în } Y,
\]


urmează că $x \in \operatorname{Dom}(T)$ și $Tx = y$.
\end{definition}

\begin{theorem}[Fréchet]
Fie $X$ și $Y$ două spații Banach și $T : X \to Y$ un operator liniar definit pe tot spațiul $X$. Atunci următoarele afirmații sunt echivalente:
\begin{enumerate}
    \item[(i)] $T$ este continuu;
    \item[(ii)] $T$ este închis;
    \item[(iii)] Graficul lui $T$,
    

\[
    \Gamma(T) = \{(x,Tx) : x \in X\},
    \]


    este închis în $X \times Y$.
\end{enumerate}
\end{theorem}

\begin{proposition}
În ipotezele teoremei lui Fréchet, următoarele condiții sunt echivalente pentru un operator liniar $T : X \to Y$:
\begin{enumerate}
    \item[(i)] Există o constantă $C > 0$ astfel încât
    

\[
    \|Tx\| \le C \|x\| \quad \text{pentru toate } x \in X.
    \]


    \item[(ii)] Graficul $\Gamma(T)$ este închis.
    \item[(iii)] $T$ este continuu în $0$.
\end{enumerate}
\end{proposition}

\end{document}
